\documentclass[a4paper,11pt]{article}
\usepackage[T1]{fontenc}
\usepackage[utf8]{inputenc}
\usepackage{t1enc}
\usepackage[dvips]{color}
\usepackage{amsmath,amssymb,makeidx,stmaryrd,yhmath}
\usepackage[normalem]{ulem}
\usepackage{fancybox}
\definecolor{gris}{gray}{0.9}
\usepackage{pst-plot,pst-tree,pstricks,pst-node,pst-text}
\usepackage{pst-eucl}
\usepackage{pstricks-add}
\usepackage[left=3.5cm, right=3.5cm, top=3cm, bottom=3cm]{geometry}
\newcommand{\vect}[1]{\overrightarrow{\,\mathstrut#1\,}}
\def\bas{\fcolorbox{gris}{gris}{\parbox{0.986\textwidth}{
	\color{red}\footnotesize 1990\hfill Page : \thepage/1}}}
\def\haut{\fcolorbox{gris}{gris}{\parbox{0.986\textwidth}{
	\color{red}
	\bf Concours Général de Mathématiques \hfill 1990
	}}}
\makeatletter
\def\ps@cg{%
\let\@oddhead\haut%
\let\@evenhead\haut%
\let\@oddfoot\bas
\let\@evenfoot\bas}
\makeatother
\newcommand{\R}{\mathbb{R}}
\newcommand{\N}{\mathbb{N}}
\newcommand{\D}{\mathbb{D}}
\newcommand{\Z}{\mathbb{Z}}
\newcommand{\Q}{\mathbb{Q}}
\newcommand{\C}{\mathbb{C}}
\pagestyle{cg}
\renewcommand{\theenumi}{\textbf{\arabic{enumi}}}
\renewcommand{\labelenumi}{\textbf{\theenumi.}}
\renewcommand{\theenumii}{\textbf{\alph{enumii}}}
\renewcommand{\labelenumii}{\textbf{\theenumii.}}
\parindent0pt
%\everymath{\displaystyle}
\usepackage[french]{babel}

\DecimalMathComma
\usepackage[np]{numprint}
\begin{document}

\begin{center}\textbf{{\color{blue} Exercice I}}
\end{center}

Soit $\left(u_n\right)_{n \in \N}$ la suite de nombres entiers définie par : 

\[\left\{\begin{array}{l c l}
u_0			&=&0 \: \text{et pour tout entier naturel}\: n\\
u_{2n}		&=& u_n \: \text{et}\\
u_{2n+ 1} 	&=& 1 - u_n.
\end{array}\right.\]

\medskip

\begin{enumerate}
\item Calculer $u_{\np{1990}}$.
\item Déterminer le nombre d'indices $n$, inférieurs ou égaux à \np{1990}, tels que $u_n = 0$.
\item Soit $p$ un nombre entier naturel et $N = (2p - 1)^2$. 

Calculer $u_N$.
\end{enumerate}

\begin{center}\textbf{{\color{blue} Exercice II}}
\end{center}

Un jeu est constitué de pièces en forme de tétraèdres réguliers d'arêtes de longueur 1.

Toutes ces pièces ont été peintes à l'aide d'une palette de $n$ couleurs : chaque face d'un tétraèdre est peinte d'une seule couleur et on précise que les quatre faces d'un tétraèdre ne sont pas nécessairement de couleurs distinctes. 

Déterminer le nombre maximum de pièces de ce jeu sachant que le jeu ne contient pas deux pièces identiques.

\begin{center}\textbf{{\color{blue} Exercice III}}
\end{center}

\begin{enumerate}
\item Trouver trois nombres entiers naturels $a, b, c$, distincts ou non, tels que:

\[\dfrac14 = \dfrac1{a^2}  + \dfrac1{b^2} + \dfrac1{c^2}\]

\item Déterminer tous les nombres entiers naturels $n$ tels qu'il existe $n$ nombres entiers naturels $x_1, x_2, ... , x_n$ , distincts ou non, vérifiant :


\[1 = \dfrac1{x_1^2}  + \dfrac1{x_2^2} + \ldots+ \dfrac1{x_n^2}.\]

\end{enumerate}

\begin{center}\textbf{{\color{blue} Exercice IV}}
\end{center}

\begin{enumerate}
\item Quelle est l'aire maximale d'un triangle dont les sommets sont dans un carré donné?
\item Quelle est le volume maximal d'un tétraèdre dont les sommets sont dans un cube donné?
\end{enumerate}

\begin{center}\textbf{{\color{blue} Exercice V}}
\end{center}

Pour tout triangle ABC, on note :
\begin{itemize}
\item[$\bullet~~$]$\Gamma$ son cercle exinscrit dans l'angle du triangle de sommet A ;
\item[$\bullet~~$]A$'$, B$'$, C$'$ les points de contact respectifs de $\Gamma$ avec les droites (BC), (CA) et (AB) ; 
$\Gamma$$S$(ABC) l'aire de la portion du plan délimitée par les segments [AB$'$], [AC$'$] et l'arc $\wideparen{\text{C}'\text{A}'\text{B}'}$ de $\Gamma$.
\end{itemize}

\begin{center}
\psset{unit=1cm}
\begin{pspicture}(10,8)
%\psgrid
\pscircle(5.7,3.8){2.9}
\psline(8,0.58)(0.2,1.6)(5.4,8)

\uput[l](0.2,1.6){A}\uput[ul](2.8,4.8){B} \uput[d](2.8,1.24){C} 
\uput[ur](7.6,6){$\Gamma$}
\uput[d](5.4,0.92){B$'$}\uput[ul](3.42,5.57){C$'$}\uput[r](2.8,3.7){A$'$}
\pscustom[fillstyle=solid,fillcolor=lightgray]
{
\psline(5.4,0.92)(0.2,1.6)(3.42,5.57)
\psarc(5.7,3.8){2.9}{133}{264}
}
\psline(2.8,4.8)(2.8,1.24)
\psdots[dotscale=0.8](0.2,1.6)(2.8,4.8)(2.8,1.24)(2.8,3.7)(5.4,0.92)(3.42,5.57)
\end{pspicture}
\end{center}

Montrer qu'il existe des triangles ABC, de périmètre $2p$ donné, pour lesquels l'aire $S$(ABC) est maximale. 

Pour un tel triangle, donner une valeur approchée, à un degré
près, de la mesure de l'angle du triangle ABC de sommet A.
\end{document}